\documentclass{article}
\usepackage{amsmath}
\usepackage{amssymb}
\usepackage{fullpage}
\usepackage{enumerate}
\usepackage{xcolor}
\usepackage{bm}
\pagestyle{empty}
\usepackage{graphicx}
\graphicspath{{C:\Users\steve\Pictures}}
\begin{document}

\noindent{{\Large \bf Fall 2025, 4100 Problem Set 2}

\large

\bigskip
\noindent{Please complete all problems. For any proof questions, please structure your proof similar to those from lecture.

\bigskip

\begin{enumerate}

\item Given $A=\left[
\begin{array}{cc}
3 & -1\\
0 & 2
\end{array}\right]$ show that the set $C=\{B\in\mathbb{R}^{2\times 2}:AB=BA\}$ is a subspaces of $\mathbb{R}^{2\times 2}$ and find a basis.


\item Let $V$ be a vector space such that $dimV=n$. Show that any set of $n$ (non-zero) vectors, $\{{\bm v}_1,\dots,{\bm v}_n\}$, satisfying ${\bm v}_i^T{\bm v}_j=0$ for $i\neq j$ is a basis for $V$.

\item  Let $\mathbb{P}_n$ be the space of all polynomials over $\mathbb{R}$ of degree at most $n$. Use the Rank-Nullity theorem to show that the subset $S=\{g\in \mathbb{P}_n: f'+4f=g~\textrm{for some}~f\in \mathbb{P}_n\}$ is a subspace of $\mathbb{P}_n$ of dimension $n+1$.

\item Let $A\in\mathbb{R}^{m\times n}$ and $B\in\mathbb{R}^{n\times l}$ such that $AB{\bm x}=\textbf{0}$ for all ${\bm x}\in \mathbb{R}^l$. Prove that ${\rm rank}(A)+{\rm rank}(B)\leq n$ using the Rank-Nullity theorem.

\item Let $V$ be a vector space with basis $\{{\bm v}_1,\cdots,{\bm v}_{n-1}\}$ and let $T\in\mathcal{L}(V,V)$ be defined by $T{\bm v}_i={\bm v}_{n-i}$. Show that $T^{-1}=T$ and that $T-I$ is not invetible. (Here $(T-I){\bm x}=T{\bm x}-{\bm x}$)

\item Consider the block-diagonal matrix $A=\left[
\begin{array}{cc}
B & 0\\
0 & C
\end{array}
\right]$. Show that ${\rm nullity}(A)={\rm nullity}(B)+{\rm nullity}(C)$.

\end{enumerate}

\end{document}